\paragraph{Predictive usefulness is a diagnostic stack.}
The results above suggest that world-model predictions should not be evaluated by a single scalar accuracy. In SPSM, predictive usefulness decomposes into five levels. First, exact-intervention hard negatives test whether a predictor is counterfactually discriminative. Second, same-state/different-action and same-action/different-state negatives localize failures to action or state axes. Third, the compute-regime map separates queries where extra predictive compute is unnecessary, useful, insufficient, or harmful. Fourth, the predicted-latent calibration map shows that latent predictions can contain physical information without being directly decodable by probes trained on true future latents. Fifth, downstream decision diagnostics test whether calibrated predicted latents support action choice. This stack is the main methodological point of SPSM: it turns latent world-model evaluation from average prediction quality into an interventional diagnostic of reliability, compute value, calibration, and actionability.
